\documentclass[10pt,a4paper]{amsart}
\usepackage{amsmath,amssymb,amsthm,graphicx}
\usepackage{verbatim}
\newenvironment{reference}{\comment}{\endcomment}
%\newenvironment{reference}{}{}
\newenvironment{slogan}{\comment}{\endcomment}
\newenvironment{history}{\comment}{\endcomment}

% For commutative diagrams we use Xy-pic
\usepackage[all]{xy}

% We use 2cell for 2-commutative diagrams.
\xyoption{2cell}
% Theorem environments.
%
\theoremstyle{plain}
\newtheorem{theorem}[subsection]{Theorem}
\newtheorem{proposition}[subsection]{Proposition}
\newtheorem{lemma}[subsection]{Lemma}

\theoremstyle{definition}
\newtheorem{definition}[subsection]{Definition}
\newtheorem{example}[subsection]{Example}
\newtheorem{exercise}[subsection]{Exercise}
\newtheorem{situation}[subsection]{Situation}

\theoremstyle{remark}
\newtheorem{remark}[subsection]{Remark}
\newtheorem{remarks}[subsection]{Remarks}

\numberwithin{equation}{subsection}

% Macros
%
\def\lim{\mathop{\mathrm{lim}}\nolimits}
\def\colim{\mathop{\mathrm{colim}}\nolimits}
\def\Spec{\mathop{\mathrm{Spec}}}
\def\Hom{\mathop{\mathrm{Hom}}\nolimits}
\def\Ext{\mathop{\mathrm{Ext}}\nolimits}
\def\SheafHom{\mathop{\mathcal{H}\!\mathit{om}}\nolimits}
\def\SheafExt{\mathop{\mathcal{E}\!\mathit{xt}}\nolimits}
\def\Sch{\mathit{Sch}}
\def\Mor{\mathop{\mathrm{Mor}}\nolimits}
\def\Ob{\mathop{\mathrm{Ob}}\nolimits}
\def\Sh{\mathop{\mathit{Sh}}\nolimits}
\def\NL{\mathop{N\!L}\nolimits}
\def\CH{\mathop{\mathrm{CH}}\nolimits}
\def\proetale{{pro\text{-}\acute{e}tale}}
\def\etale{{\acute{e}tale}}
\def\QCoh{\mathit{QCoh}}
\def\Ker{\mathop{\mathrm{Ker}}}
\def\Im{\mathop{\mathrm{Im}}}
\def\Coker{\mathop{\mathrm{Coker}}}
\def\Coim{\mathop{\mathrm{Coim}}}

% Boxtimes
%
\DeclareMathSymbol{\boxtimes}{\mathbin}{AMSa}{"02}

%
% Macros for moduli stacks/spaces
%
\def\QCohstack{\mathcal{QC}\!\mathit{oh}}
\def\Cohstack{\mathcal{C}\!\mathit{oh}}
\def\Spacesstack{\mathcal{S}\!\mathit{paces}}
\def\Quotfunctor{\mathrm{Quot}}
\def\Hilbfunctor{\mathrm{Hilb}}
\def\Curvesstack{\mathcal{C}\!\mathit{urves}}
\def\Polarizedstack{\mathcal{P}\!\mathit{olarized}}
\def\Complexesstack{\mathcal{C}\!\mathit{omplexes}}
% \Pic is the operator that assigns to X its picard group, usage \Pic(X)
% \Picardstack_{X/B} denotes the Picard stack of X over B
% \Picardfunctor_{X/B} denotes the Picard functor of X over B
\def\Pic{\mathop{\mathrm{Pic}}\nolimits}
\def\Picardstack{\mathcal{P}\!\mathit{ic}}
\def\Picardfunctor{\mathrm{Pic}}
\def\Deformationcategory{\mathcal{D}\!\mathit{ef}}
\usepackage{cleveref}
\crefname{originallemma}{original lemma}{original lemmas}
\Crefname{originallemma}{Original lemma}{Original lemmas}
\makeatletter
\newlabel{more-algebra-lemma-base-change-comparison}{{0661}{source}{}{sourcepointer-1}{}}
\newlabel{more-algebra-lemma-base-change-comparison@cref}{{[originallemma][0][]0661}{[1][0][]source}{}{}{}}
\makeatother
\crefname{originallemma}{original lemma}{original lemmas}
\Crefname{originallemma}{Original lemma}{Original lemmas}
\makeatletter
\newlabel{more-algebra-lemma-pull-tor-amplitude}{{066L}{source}{}{sourcepointer-2}{}}
\newlabel{more-algebra-lemma-pull-tor-amplitude@cref}{{[originallemma][0][]066L}{[1][0][]source}{}{}{}}
\makeatother
\crefname{originallemma}{original lemma}{original lemmas}
\Crefname{originallemma}{Original lemma}{Original lemmas}
\makeatletter
\newlabel{sites-cohomology-lemma-tor-amplitude-stalk}{{0DJJ}{source}{}{sourcepointer-3}{}}
\newlabel{sites-cohomology-lemma-tor-amplitude-stalk@cref}{{[originallemma][0][]0DJJ}{[1][0][]source}{}{}{}}
\makeatother
\crefname{originaltheorem}{original theorem}{original theorems}
\Crefname{originaltheorem}{Original theorem}{Original theorems}
\makeatletter
\newlabel{spaces-properties-theorem-exactness-stalks}{{04K5}{source}{}{sourcepointer-4}{}}
\newlabel{spaces-properties-theorem-exactness-stalks@cref}{{[originaltheorem][0][]04K5}{[1][0][]source}{}{}{}}
\makeatother
\title{Tor independence preserves tor amplitude under base change}
\author{The Stacks Project Authors}
\date{stacks-project:spaces-perfect.tex:local-snapshot-20261012; GNU FDL 1.2 or later}
\begin{document}
\maketitle
\noindent\textit{Editorial note.} All delivered mathematical text is the author's own, copied from the pinned Stacks Project chapter file spaces-perfect.tex; no statement or proof step was rewritten, reordered or removed.\par\medskip
\noindent\textit{Editorial note.} The counted proof is the author's complete argument: the proof computes the derived pullback on stalks, applies the tor-independence base change comparison, derives part (1) from the pull-tor-amplitude lemma, and reads part (2) off part (1) via the tor amplitude characterization of flatness\par\medskip
\noindent\textit{Editorial note.} The label lemma-tor-independence-and-tor-amplitude names the counted statement, and the proof environment immediately following it in the source is delivered verbatim as its proof.\par\medskip
\noindent\textit{Editorial note.} The counted statement is carried by the author's own lemma environment in the chapter source.\par\medskip
\noindent\textit{Allowed standard prerequisites (editorial).}\begin{itemize}
\item tor amplitude of complexes computed on stalks (Cohomology on Sites, Lemma sites-cohomology-lemma-tor-amplitude-stalk) with the stalk theorem for exactness (Properties of Spaces, Theorem spaces-properties-theorem-exactness-stalks)
\item the stalk formula for derived pullback along the base changed morphism
\item the base change comparison of derived tensor products under tor independence (More on Algebra, Lemma more-algebra-lemma-base-change-comparison)
\item the behaviour of tor amplitude under field/tensor base change (More on Algebra, Lemma more-algebra-lemma-pull-tor-amplitude)
\end{itemize}\medskip
\ifdefined\hypertarget\hypertarget{sourcepointer-1}{}\fi
\noindent\textit{Source reference.} Original lemma 0661: sources/stacks-chow/tags.txt, source lines 7920--7920. Prerequisite pointer invoked at proof line(s) 5269; it names a standard fact of a sibling chapter that the counted argument uses, and every step of the counted proof built on it is retained verbatim in the delivered body.\par\medskip
\ifdefined\hypertarget\hypertarget{sourcepointer-2}{}\fi
\noindent\textit{Source reference.} Original lemma 066L: sources/stacks-chow/tags.txt, source lines 7940--7940. Prerequisite pointer invoked at proof line(s) 5276; it names a standard fact of a sibling chapter that the counted argument uses, and every step of the counted proof built on it is retained verbatim in the delivered body.\par\medskip
\ifdefined\hypertarget\hypertarget{sourcepointer-3}{}\fi
\noindent\textit{Source reference.} Original lemma 0DJJ: sources/stacks-chow/tags.txt, source lines 16979--16979. Prerequisite pointer invoked at proof line(s) 5255; it names a standard fact of a sibling chapter that the counted argument uses, and every step of the counted proof built on it is retained verbatim in the delivered body.\par\medskip
\ifdefined\hypertarget\hypertarget{sourcepointer-4}{}\fi
\noindent\textit{Source reference.} Original theorem 04K5: sources/stacks-chow/tags.txt, source lines 5961--5961. Prerequisite pointer invoked at proof line(s) 5257; it names a standard fact of a sibling chapter that the counted argument uses, and every step of the counted proof built on it is retained verbatim in the delivered body.\par\medskip
\begin{lemma}
\label{lemma-tor-independence-and-tor-amplitude}
Let $S$ be a scheme. Consider a cartesian square of algebraic spaces
$$
\xymatrix{
X' \ar[r]_{g'} \ar[d]_{f'} & X \ar[d]^f \\
Y' \ar[r]^g & Y
}
$$
over $S$. Assume $g$ and $f$ Tor independent.
\begin{enumerate}
\item If $E \in D(\mathcal{O}_X)$ has tor amplitude
in $[a, b]$ as a complex of $f^{-1}\mathcal{O}_Y$-modules,
then $L(g')^*E$ has tor amplitude
in $[a, b]$ as a complex of $f^{-1}\mathcal{O}_{Y'}$-modules.
\item If $\mathcal{G}$ is an $\mathcal{O}_X$-module flat
over $Y$, then $L(g')^*\mathcal{G} = (g')^*\mathcal{G}$.
\end{enumerate}
\end{lemma}

\begin{proof}
We can compute tor dimension at stalks, see
Cohomology on Sites, Lemma \ref{sites-cohomology-lemma-tor-amplitude-stalk}
and Properties of Spaces, Theorem
\ref{spaces-properties-theorem-exactness-stalks}.
If $\overline{x}'$ is a geometric point of $X'$ with image
$\overline{x}$ in $X$, then
$$
(L(g')^*E)_{\overline{x}'} =
E_{\overline{x}}
\otimes_{\mathcal{O}_{X, \overline{x}}}^\mathbf{L}
\mathcal{O}_{X', \overline{x}'}
$$
Let $\overline{y}'$ in $Y'$ and $\overline{y}$ in $Y$
be the image of $\overline{x}'$ and $\overline{x}$.
Since $X$ and $Y'$ are tor independent over $Y$, we can apply
More on Algebra, Lemma \ref{more-algebra-lemma-base-change-comparison}
to see that the right hand side of the displayed formula is equal to
$E_{\overline{x}}
\otimes_{\mathcal{O}_{Y, \overline{y}}}^\mathbf{L}
\mathcal{O}_{Y', \overline{y}'}$
in $D(\mathcal{O}_{Y', \overline{y}'})$.
Thus (1) follows from
More on Algebra, Lemma \ref{more-algebra-lemma-pull-tor-amplitude}.
To see (2) observe that flatness of $\mathcal{G}$ is equivalent to
the condition that $\mathcal{G}[0]$ has tor amplitude in $[0, 0]$.
Applying (1) we conclude.
\end{proof}

\end{document}
